\section{Del salto espectral \texorpdfstring{\(4\to2\to6\to54\)}{4 a 2 a 6 a 54}: acoplamiento local, compresión modular e integración bidimensional}
\label{sec:sucesor83-salto-espectral-4-2-6-54}

Para dos representantes consecutivos \((k,k+1)\), considérese
\[
B_k=
\begin{pmatrix}
k^2&k(k+1)\\
k(k+1)&(k+1)^2
\end{pmatrix}\pmod 9.
\]

\begin{theorem}[Unicidad del bloque diagonal adyacente]
\label{thm:sucesor83-app-bloque-central}
Existe un único bloque \(B_k\) cuyas entradas diagonales coinciden. Está
determinado por \(k=4\) y es
\[
B_c=\begin{pmatrix}7&2\\2&7\end{pmatrix}.
\]
\end{theorem}

\begin{proof}
La igualdad de las diagonales equivale a
\[
k^2\equiv(k+1)^2\pmod 9,
\]
es decir, \(2k+1\equiv0\pmod 9\). Como \(2\) es una unidad de
\(\mathbb Z/9\mathbb Z\), la congruencia posee la única solución
\(k\equiv4\pmod 9\).
\end{proof}

Por tanto,
\[
B_c=7I+2R,
\qquad
R=\begin{pmatrix}0&1\\1&0\end{pmatrix},
\qquad R^2=I.
\]
Con \(P_\pm=(I\pm R)/2\),
\[
B_c=9P_++5P_-.
\]

\begin{theorem}[Jerarquía de la separación entre suma y producto]
\label{thm:sucesor83-app-jerarquia-4-2-6-54}
El bloque central tiene salto espectral local \(9-5=4\). Su coeficiente
fuera de la diagonal es \(2\). La compresión por las \(3\) clases módulo
\(3\) produce la diferencia agregada \(51-45=6\), y el despliegue sobre las
\(9\) posiciones produce el defecto global
\[
9\cdot6=459-405=54.
\]
Los enteros \(4,2,6,54\) pertenecen, respectivamente, al salto espectral,
al acoplamiento local, a la compresión modular y a la integración
bidimensional.
\end{theorem}

\begin{proof}
La descomposición \(B_c=7I+2R=9P_++5P_-\) da el salto y el coeficiente
local. Las matrices agregadas \(M_\Pi,M_\Sigma\) satisfacen
\(\sum M_\Pi=51\) y \(\sum M_\Sigma=45\). Cada entrada agregada representa
\(9\) posiciones, por lo que el despliegue multiplica la diferencia por
\(9\) y recupera los totales \(459\) y \(405\).
\end{proof}

Sus concatenaciones ordenadas producen los bloques \(72\) y \(27\):
\[
72+27=99,
\qquad
72-27=45=\det B_c.
\]
La concatenación de las entradas diagonales y antidiagonales produce,
respectivamente, \(77\) y \(22\). Por tanto,
\[
72+27=77+22=99,
\qquad
77-22=55=54+1.
\]
Estas igualdades expresan dos particiones exactas del mismo bloque; no
identifican por sí solas una constante externa.
